\documentclass[11pt]{article}
\usepackage[T1]{fontenc}
\usepackage{lmodern}
\usepackage[margin=1in]{geometry}
\usepackage{amsmath,amssymb,amsthm,mathtools}
\usepackage{booktabs,tabularx,array}
\usepackage[round,authoryear]{natbib}
\usepackage{microtype}
\usepackage[colorlinks=true,linkcolor=blue,citecolor=blue,urlcolor=blue]{hyperref}
\newtheorem{theorem}{Theorem}[section]
\newtheorem{proposition}[theorem]{Proposition}
\newtheorem{corollary}[theorem]{Corollary}
\newtheorem{remark}[theorem]{Remark}
\newcommand{\KL}{\operatorname{KL}}
\newcommand{\Tr}{\operatorname{tr}}
\newcommand{\E}{\mathbb E}
\newcommand{\R}{\mathbb R}
\newcommand{\SPD}{\operatorname{Sym}_{++}}
\newcommand{\GL}{\operatorname{GL}}
\newcommand{\Ga}{\operatorname{Gamma}}
\newcommand{\diag}{\operatorname{diag}}
\newcommand{\Cov}{\operatorname{Cov}}
\newcommand{\Var}{\operatorname{Var}}
\newcommand{\m}{\boldsymbol{\mu}}
\newcommand{\Sig}{\boldsymbol{\Sigma}}
\setlength{\emergencystretch}{3em}
\title{Nonlinear Inference and Information Control in Gauge-VFE Agent Networks\\[0.5em]\large A Conditional Extension of Gauge-Covariant Gaussian Attention}
\author{Robert C. Dennis\\Independent Researcher}
\date{September 20, 2026}
\begin{document}
\maketitle
\begin{abstract}
The gauge-covariant Gaussian attention construction leaves the coordinate MLP, normalization, and full-stack inference interpretation unresolved. We develop a conditional extension in which interacting agents retain full Gaussian recognition states, while finite source labels and latent message precisions supply nonlinear processing. A normalized population model yields exact categorical, Gamma, and Gaussian coordinate updates. Constraining the Gaussian coordinate by a reference-relative information budget gives a unique exponential-family interpolation that preserves positive-definite covariance, controls relative belief scale, and decreases the same variational free energy under sequential updates. A two-agent example establishes a nonaffine response with cross-agent covariance and a finite budget. We also derive an independent Gaussian Fisher-flow construction, a latent-Gaussian categorical observation model, and counterexamples to universal identifications with LayerNorm or arbitrary MLPs. The results establish compatible statistical mechanisms without prescribed neuron layers. They do not select a unique architecture, reproduce the existing KL-attention score in general, or establish efficient causal language modeling. Priors, observation laws, information budgets, computational schedules, and scalable approximation remain explicit modeling obligations.
\end{abstract}

\section{Scope and continuity with the Gaussian attention construction}
\label{sec:scope}

The companion manuscripts \path{GL(K)_attention.tex} and \path{GL(K)_supplementary.tex} formulate full Gaussian site states, invariant comparison scores, and independent predictive communication. Their matrix-link construction identifies typed sending and receiving maps whose compositions transform covariantly under independent site changes of frame. Their layer-composition discussion also states the unresolved boundary: the implemented coordinate LayerNorm and residual MLP are not derived from the attention objective, and covariance passthrough does not propagate the uncertainty of the nonlinear mean map. This extension addresses that boundary by constructing complete probabilistic alternatives for several computational roles.

An agent network is retained throughout. The objects being removed are prescribed neuron layers and a fixed block that orders attention, normalization, and an MLP. Agents may communicate through attention networks, singular matrix links, and multiple channels. Their computational response may be nonlinear even when every declared state-to-record map is linear: inference over uncertain source assignment and message precision already produces nonlinear functions of observations and beliefs. The resulting theory remains an extension of the gauge-VFE program because its states, links, passive transformation laws, and variational objective have the same geometric roles. The additional statistical laws are specified assumptions, rather than consequences of the frame symmetry alone.

The main construction in Sections~\ref{sec:population}--\ref{sec:gauge} combines nonlinear inference and information control in one constrained negative evidence lower bound, or negative ELBO. Section~\ref{sec:geometry} derives its frozen Gaussian-block dynamics. Section~\ref{sec:globalflow} presents a separate energy-based construction with a global flow theorem; its hypotheses are not transferred to the first model. Sections~\ref{sec:readout} and~\ref{sec:causality} treat categorical observations and autoregression. Exact component obstructions and remaining choices are stated in Sections~\ref{sec:obstructions} and~\ref{sec:choices}.

These mechanisms have statistical precedents. Coordinate variational inference is developed by \citet{Blei2017}; Gaussian scale mixtures by \citet{AndrewsMallows1974}; and latent-Gaussian categorical augmentation by \citet{AlbertChib1993}. Objective-based transformer accounts also exist, including CRATE's compression and sparsity construction \citep{Yu2024CRATE}. The contribution here is the explicit compatibility of declared inference mechanisms with full Gaussian agent states, typed links, and reference-relative scale control. No priority claim is made for source mixtures, latent precision, or natural-gradient descent.

\section{Agents, links, and a normalized population law}
\label{sec:population}

Fix one base point and a finite population of $N\ge2$ agents, each with state space $E_i\simeq\R^K$, $K\ge2$. Collect the latent states into $x=(x_1,\ldots,x_N)\in\R^D$, where $D=NK$. The recognition state will be a full population Gaussian,
\begin{equation}
Q(x)=\mathcal N(x;\m,\Sig), \Sig\in\SPD(D).
\label{eq:jointgaussian}
\end{equation}
Here $\m_i$ and $\Sig_{ii}$ are the mean and full covariance of agent $i$, while $\Sig_{ij}$ represents cross-agent covariance. This notation distinguishes the population covariance from the reference-coordinate covariances denoted $C_i$ in the companion paper. A product of site Gaussians is a further restriction, not the default object in this extension.

The companion matrix links can be retained as
\begin{equation}
L_{ij}^{(h)}=U_i O_h T_h U_j^{-1},
L_{ij}^{(h)\prime}=g_iL_{ij}^{(h)}g_j^{-1}.
\label{eq:inheritedlinks}
\end{equation}
Independent typed links obeying the same endpoint law are also admissible. A singular link is a communication map; it is not automatically an invertible connection. The general-matrix link vocabulary has a lattice-gauge precedent \citep{Wetterich2014}, but no lattice action, continuum limit, or compact-group conclusion is imported into the present $\GL(K)$ model.

Declare a proper reference law
\begin{equation}
r(x)=\mathcal N(x;\m_0,\mathsf P^{-1}), \mathsf P\succ0.
\label{eq:reference}
\end{equation}
Let $a$ index finitely many observation or interaction records. Record $a$ has dimension $d_a\ge1$, a finite nonempty source palette $\mathcal J_a$, positive probabilities $\pi_{aj}$ summing to one, receiving noise $R_a\in\SPD(d_a)$, and a positive parameter $\nu_a$. For every $j\in\mathcal J_a$, specify a linear map $\mathsf A_{aj}:\R^D\to\R^{d_a}$ and offset $b_{aj}\in\R^{d_a}$. One realization connects these maps directly to the existing links:
\begin{equation}
\mathsf A_{aj}x=\mathsf H_a(x_i-L_{ij}^{(h)}x_j),
\label{eq:recordlink}
\end{equation}
where $a$ belongs to receiver $i$, $\mathsf H_a:E_i\to\R^{d_a}$ is a declared measurement, and a palette may enumerate sources or source-channel pairs. Neither $\mathsf A_{aj}$ nor $L_{ij}^{(h)}$ must be invertible. Receiving noise remains positive definite even when the measured or communicated subspace has lower rank.

Introduce a source label $z_a$ and positive reliability $\lambda_a$. With Gamma densities parameterized by shape and rate, define
\begin{equation}
\begin{split}
p(x,z,\lambda,o)
={}&r(x)\prod_a \pi_{a z_a}
\Ga(\lambda_a;\nu_a/2,\nu_a/2)\\
&\mathrel{}\times\mathcal N(o_a;\mathsf A_{a z_a}x+b_{a z_a},R_a/\lambda_a).
\end{split}
\label{eq:jointlaw}
\end{equation}
The base measures are Lebesgue measure on $x$ and $o$, positive Lebesgue measure on $\lambda$, and counting measure on $z$. All model parameters and observations are finite. The reference, palettes, maps, noise matrices, and reliability laws are model choices. No unspecified neural map is hidden in the linear record operators.

\begin{proposition}[Normalization and finite evidence]
\label{prop:normalization}
Equation~\eqref{eq:jointlaw} defines a probability law. For every finite observed record vector $o$, its evidence satisfies $0<p(o)<\infty$.
\end{proposition}
\begin{proof}
Each Gaussian record integrates to one conditional on $(x,z,\lambda)$. Tonelli's theorem applies to the nonnegative joint density. Integrating records, summing labels, integrating the proper Gamma factors, and integrating $r$ gives total mass one. This argument permits cycles among the agent pairs appearing in \eqref{eq:recordlink}; all records are generated from a common latent population. It does not multiply mutually recursive state conditionals.

For fixed $x$ and $z_a=j$, write $e_{aj}=o_a-b_{aj}-\mathsf A_{aj}x$ and $s_{aj}=e_{aj}^{\top}R_a^{-1}e_{aj}$. Direct Gamma integration gives
\begin{equation}
p(o_a\mid x,z_a=j)=
\frac{\Gamma((\nu_a+d_a)/2)}{\Gamma(\nu_a/2)(\pi\nu_a)^{d_a/2}|R_a|^{1/2}}
\left(1+\frac{s_{aj}}{\nu_a}\right)^{-(\nu_a+d_a)/2}.
\label{eq:student}
\end{equation}
This density is positive everywhere and bounded above by its finite leading constant. A finite product of the corresponding finite mixtures is positive and bounded. Integrating against the proper reference proves the evidence claim. No finite-variance assumption on the Student law is needed; $\nu_a>0$ suffices.
\end{proof}

\section{One variational objective and its exact coordinates}
\label{sec:coordinates}

Use the recognition family
\begin{equation}
q(x,z,\lambda)=Q(x)\prod_a\beta_a(z_a)\gamma_a(\lambda_a),
\label{eq:recognition}
\end{equation}
where $Q$ is \eqref{eq:jointgaussian}. The auxiliary factors are independent of the population state under recognition; the true posterior generally has neither that independence nor Gaussian state marginals. Initially require finite negative ELBO, finite $\E_{\gamma_a}\lambda_a$, and finite $\E_{\gamma_a}|\log\lambda_a|$. Positive finite-shape/rate Gamma factors satisfy these conditions. They remain of that form after their coordinate update. Mere feasibility of a Gaussian information budget would not imply these integrability conditions.

Put $w_a=\E_{\gamma_a}\lambda_a$, $\ell_a=\E_{\gamma_a}\log\lambda_a$, and
\begin{equation}
\begin{split}
\delta_{aj}={}&(o_a-b_{aj}-\mathsf A_{aj}\m)^\top
 R_a^{-1}(o_a-b_{aj}-\mathsf A_{aj}\m)\\
&+\Tr(R_a^{-1}\mathsf A_{aj}\Sig\mathsf A_{aj}^\top).
\end{split}
\label{eq:delta}
\end{equation}
The trace is part of the expected residual energy. For \eqref{eq:recordlink}, its unmeasured covariance contribution contains
\begin{equation}
\Sig_{ii}+L_{ij}^{(h)}\Sig_{jj}L_{ij}^{(h)\top}
-\Sig_{ij}L_{ij}^{(h)\top}-L_{ij}^{(h)}\Sig_{ji}.
\label{eq:crosscov}
\end{equation}
A product approximation removes the last two terms by assumption. They cannot be omitted while claiming a full correlated Gaussian calculation.

Expanding $\E_q\log(q/p(x,z,\lambda,o))$ gives exactly
\begin{align}
\mathcal F={}&\KL(Q\Vert r)+\sum_a\KL(\beta_a\Vert\pi_a)
 +\sum_a\KL\bigl(\gamma_a\Vert\Ga(\nu_a/2,\nu_a/2)\bigr)\nonumber\\
&+\frac12\sum_a\left[d_a\log(2\pi)+\log|R_a|-d_a\ell_a
 +w_a\sum_j\beta_{aj}\delta_{aj}\right].
\label{eq:freeenergy}
\end{align}
There is one population-state entropy and one entropy for each auxiliary factor. The posterior identity is
\begin{equation}
\mathcal F=\KL\bigl(q\Vert p(x,z,\lambda\mid o)\bigr)-\log p(o).
\label{eq:elboidentity}
\end{equation}
Thus an exact negative ELBO is obtained for a restricted family; exactness of this identity does not mean that recognition equals the posterior.

\begin{proposition}[Coordinate solutions]
\label{prop:coordinates}
At fixed values of the other factors, the unique source and reliability coordinate minimizers are
\begin{align}
\beta_{aj}^*&=\frac{\pi_{aj}\exp(-w_a\delta_{aj}/2)}
 {\sum_{k\in\mathcal J_a}\pi_{ak}\exp(-w_a\delta_{ak}/2)},
\label{eq:betastar}\\
\gamma_a^*&=\Ga\left(\frac{\nu_a+d_a}{2},
 \frac{\nu_a+\sum_j\beta_{aj}\delta_{aj}}2\right).
\label{eq:gammastar}
\end{align}
The unconstrained state coordinate is $Q_* =\mathcal N(\m_*,\Sig_*)$, where
\begin{align}
\mathsf M:=\Sig_*^{-1}
 &=\mathsf P+\sum_{a,j}w_a\beta_{aj}\mathsf A_{aj}^{\top}R_a^{-1}\mathsf A_{aj},
\label{eq:precision}\\
\mathsf h:=\Sig_*^{-1}\m_*
 &=\mathsf P\m_0+\sum_{a,j}w_a\beta_{aj}\mathsf A_{aj}^{\top}R_a^{-1}(o_a-b_{aj}).
\label{eq:information}
\end{align}
\end{proposition}
\begin{proof}
The terms involving one source row are $\sum_j\beta_{aj}\log(\beta_{aj}/\pi_{aj})+(w_a/2)\sum_j\beta_{aj}\delta_{aj}$. Substituting \eqref{eq:betastar} rewrites them as $\KL(\beta_a\Vert\beta_a^*)$ plus a constant. The common receiving covariance and record dimension make the determinant and expected-log-precision contributions constant within that row. If noise depends on the label, both its quadratic form and determinant must instead enter the score.

Collecting terms that depend on $\lambda_a$ in the expected log joint gives
\begin{equation}
\left(\frac{\nu_a+d_a}{2}-1\right)\log\lambda_a
-\frac{\nu_a+\sum_j\beta_{aj}\delta_{aj}}2\lambda_a+c.
\label{eq:gammacollect}
\end{equation}
The rate is strictly positive and finite. Its normalized exponential is \eqref{eq:gammastar}; the coordinate objective is its KL plus a constant, proving optimality among admissible auxiliary densities.

The state-dependent expected log joint is $-x^\top\mathsf M x/2+\mathsf h^\top x+c$. Every added precision term in \eqref{eq:precision} is positive semidefinite, so $\mathsf M\succeq\mathsf P\succ0$. Completing the square gives \eqref{eq:precision}--\eqref{eq:information}. The entire state-coordinate objective is $\KL(Q\Vert Q_*)$ plus a constant, which proves uniqueness.
\end{proof}

After a reliability replacement,
\begin{equation}
w_a^*=\frac{\nu_a+d_a}{\nu_a+\sum_j\beta_{aj}\delta_{aj}}.
\label{eq:reliability}
\end{equation}
The inferred precision decreases as the expected residual increases. Both $w_a$ and the source probabilities depend on the current means and covariances, so repeated coordinate updates supply nonlinear response even though the record means are affine. The temperature in \eqref{eq:betastar} is fixed by the declared joint. Adding an arbitrary temperature changes the objective or the probabilistic model.

\subsection{Relation to the existing Gaussian-KL attention score}
\label{sec:scorecomparison}

The companion attention operator uses a comparison cost of the form
\begin{equation}
D_{ij}=\KL(q_i^{Q}\Vert q_j^{K}),
\beta_{ij}\propto\pi_{ij}\exp(-D_{ij}/\tau),
\label{eq:oldscore}
\end{equation}
with the two projected laws expressed in a common channel. Equation~\eqref{eq:betastar} uses an expected record likelihood cost, $w_a\delta_{aj}/2$, including the joint covariance in \eqref{eq:crosscov}. The two costs generally differ. In particular, an expected quadratic residual is not the complete KL between two Gaussian marginals, which also contains entropy and determinant contributions tied to both distributions.

The correspondence established here is the source-selection mechanism and its covariance under passive frame changes. Exact recovery of \eqref{eq:oldscore} would require an additional normalized model construction making its comparison cost the appropriate coordinate energy. Arbitrarily substituting that cost into \eqref{eq:betastar} while retaining \eqref{eq:precision}--\eqref{eq:information} would no longer be the coordinate optimization proved above. Predictive content readout in the companion manuscript is likewise distinct from updating a recognition mean in this population model.

\section{A nonlinear interacting example}
\label{sec:example}

Consider first one scalar latent contrast $v\sim\mathcal N(0,1)$ and one record $y\mid v,\lambda\sim\mathcal N(v,1/\lambda)$, with $\lambda\sim\Ga(\nu/2,\nu/2)$. A singleton source palette isolates the reliability mechanism. Starting from the prior recognition law, optimize the reliability factor and then the Gaussian factor. Since $\E(y-v)^2=y^2+1$, the resulting moments are
\begin{equation}
 w^+=\frac{\nu+1}{\nu+y^2+1},
 \mu_v^+(y)=\frac{(\nu+1)y}{y^2+2\nu+2},
 \Sigma_v^+(y)=\frac{\nu+y^2+1}{y^2+2\nu+2}.
\label{eq:scalarresponse}
\end{equation}
The derivative of the mean at zero is $1/2$, whereas $\mu_v^+(y)/y\to0$ as $|y|\to\infty$. The response is therefore not affine. Its covariance also changes with the data; replacing it by a fixed passthrough would alter this inference rule.

To place the example inside the admitted interacting population, take $N=2$, $K=2$, $r=\mathcal N(0,I_4)$, and
\begin{equation}
\mathsf A=\frac1{\sqrt2}(1,0,-1,0), v=\mathsf A x.
\label{eq:contrast}
\end{equation}
The contrast has unit prior variance. Its orthogonal complement retains its prior under this coordinate replacement, so
\begin{equation}
\m^+=\mathsf A^\top\mu_v^+,
\Sig^+=I_4+(\Sigma_v^+-1)\mathsf A^\top\mathsf A.
\label{eq:embedding}
\end{equation}
The covariance eigenvalues are $\Sigma_v^+>0$ and three ones. The first coordinates of the two agents have cross covariance $(1-\Sigma_v^+)/2>0$. Extra agents and coordinates may retain independent unit-Gaussian priors, extending this witness to every finite $N\ge2$, $K\ge2$. This establishes existence of nonlinear interacting inference, rather than nonlinearity for every possible parameter choice.

\section{Information control as an exact coordinate problem}
\label{sec:budget}

A scale constraint can be integrated into the same objective. Fix a finite budget $B\ge0$ and require
\begin{equation}
\KL(Q\Vert r)\le B.
\label{eq:budget}
\end{equation}
The reference transforms with the population state and supplies the comparison geometry. Choosing its numerical parameters, the direction of KL, and the value of $B$ is additional model or inference design. This constraint controls relative belief displacement; it does not define an absolute Euclidean scale invariant under arbitrary $\GL(K)$ changes.

\begin{theorem}[Exact Gaussian reference-budget projection]
\label{thm:projection}
Let $q$ and $r$ be full nondegenerate Gaussian laws in finite dimension and $B\ge0$ finite. The unique minimizer of $\KL(s\Vert q)$ over all probability densities satisfying $\KL(s\Vert r)\le B$ is $q$ if it is feasible and $r$ if $B=0$. Otherwise it is
\begin{equation}
q_t=\frac{q^t r^{1-t}}{\int q^t r^{1-t}}, 0<t<1,
\label{eq:projection}
\end{equation}
where $t$ is the unique solution of $\KL(q_t\Vert r)=B$.
\end{theorem}
\begin{proof}
Write $L=\log(q/r)$, $A(t)=\log\E_r e^{tL}$, and $q_t=r e^{tL-A(t)}$. Gaussian precision interpolation is positive definite on $[0,1]$ and a neighborhood of that interval. Polynomial derivatives of the quadratic log-density ratio therefore have integrable Gaussian domination. Differentiating gives
\begin{equation}
A'(t)=\E_{q_t}L, A''(t)=\Var_{q_t}L,
\frac{d}{dt}\KL(q_t\Vert r)=t\Var_{q_t}L.
\label{eq:monotonicity}
\end{equation}
If $q\ne r$, the polynomial $L$ is nonconstant and a full-support Gaussian cannot concentrate on one of its level sets. The derivative is strictly positive for $t>0$. Continuity and the endpoint divergences $0$ and $\KL(q\Vert r)$ prove existence and uniqueness of the active parameter.

For any feasible competitor $s$, its second moment is finite. Indeed, for sufficiently small $a>0$, $\E_r e^{a\|x\|^2}<\infty$. The entropy inequality for truncated versions of $a\|x\|^2$ and monotone convergence give
\begin{equation}
a\E_s\|x\|^2\le\KL(s\Vert r)+\log\E_r e^{a\|x\|^2}<\infty.
\label{eq:entropybound}
\end{equation}
The entropy inequality follows by Jensen's inequality applied to $e^f r/s$ under $s$. Thus the Gaussian log-density ratios used below are integrable. Put $\eta=(1-t)/t$. Expanding logarithms and subtracting the identity at $q_t$ yields
\begin{equation}
\begin{split}
\KL(s\Vert q)-\KL(q_t\Vert q)
={}&(1+\eta)\KL(s\Vert q_t)\\
&+\eta\bigl[B-\KL(s\Vert r)\bigr]\ge0.
\end{split}
\label{eq:projectionproof}
\end{equation}
Equality forces $s=q_t$ almost everywhere. When $B=0$, KL nonnegativity makes $r$ the unique feasible law. When $q$ is feasible, its zero objective is uniquely minimal. These arguments also cover $q=r$, without asserting uniqueness of an unnecessary path parameter.
\end{proof}

Apply the theorem to the Gaussian target $Q_*$ of Proposition~\ref{prop:coordinates}. If it is infeasible, the exact constrained coordinate has
\begin{align}
\Sig_t^{-1}&=t\mathsf M+(1-t)\mathsf P,\label{eq:budgetprecision}\\
\Sig_t^{-1}\m_t&=t\mathsf h+(1-t)\mathsf P\m_0.
\label{eq:budgetmean}
\end{align}
The optimizer over all densities is itself Gaussian. Thus the Gaussian family introduces no additional loss within this particular frozen coordinate problem. A numerical implementation must still choose a tolerance or certified feasibility rule for the monotone scalar equation; a floating-point root is not an exact boundary certificate.

\begin{corollary}[Common-objective descent]
\label{cor:descent}
With fixed model data, fixed $r,B$, finite initial objective and moments, and a feasible Gaussian initialization, sequential exact replacements of $\beta$, $\gamma$, and the constrained Gaussian coordinate never increase \eqref{eq:freeenergy}. The objective values converge.
\end{corollary}
\begin{proof}
Each replacement minimizes the same objective with the other factors fixed. The old factor remains an admissible competitor. Equation~\eqref{eq:elboidentity} and Proposition~\ref{prop:normalization} give the finite lower bound $-\log p(o)$, so the decreasing sequence of objective values has a finite limit.
\end{proof}

This conclusion does not establish convergence of the parameters, a global optimum, a universal stopping depth, or descent under simultaneous stale-factor updates. Changing model parameters, the reference, or the budget during the sweep changes the optimization problem. Applying a normalization projection after an unrelated update also lacks this proof: the projection must be the constrained minimizer of the actual coordinate objective.

\subsection{Scale and marginal control}

For a feasible Gaussian $Q$, let $s_b$ be the eigenvalues of $\mathsf P^{1/2}\Sig\mathsf P^{1/2}$. The Gaussian KL identity gives
\begin{equation}
2\KL(Q\Vert r)=(\m-\m_0)^\top\mathsf P(\m-\m_0)
+\sum_{b=1}^D(s_b-\log s_b-1).
\label{eq:spectralbound}
\end{equation}
Every term is nonnegative, and $s-\log s-1$ diverges as $s\downarrow0$ or $s\to\infty$. A finite budget therefore bounds the relative mean and every relative covariance eigenvalue away from both boundaries. The KL chain rule also gives
\begin{equation}
\KL(Q_i\Vert r_i)\le\KL(Q\Vert r)\le B,
\label{eq:marginalbudget}
\end{equation}
without requiring independence of agent states.

The nonlinear witness can satisfy a fixed budget uniformly in its scalar datum. At $\nu=1$, \eqref{eq:scalarresponse} has $(\mu_v^+)^2\le1/4$ and $1/2\le\Sigma_v^+<1$. The unchanged orthogonal directions contribute zero KL, so
\begin{equation}
2\KL(Q^+\Vert r)
\le\frac14+\log2-\frac12<1.
\label{eq:witnessbudget}
\end{equation}
Hence $B=1$ preserves this nonaffine interacting response for all finite $y$. At $B=0$, by contrast, the Gaussian coordinate is fixed to $r$ and cannot provide data-dependent adaptation.
\section{Passive gauge covariance and variational restrictions}
\label{sec:gauge}

Let $G=\diag(g_1,\ldots,g_N)$ with every $g_i\in\GL(K)$. For a fixed passive coordinate change, transform
\begin{align}
x'&=Gx,& \m'&=G\m,& \Sig'&=G\Sig G^\top,\nonumber\\
\m_0'&=G\m_0,& \mathsf P'&=G^{-\top}\mathsf P G^{-1},&
\mathsf A_{aj}'&=\mathsf A_{aj}G^{-1}.
\label{eq:gaugedata}
\end{align}
Record coordinates and $o_a,b_{aj},R_a,\nu_a,\pi_{aj}$ remain fixed. For \eqref{eq:recordlink}, the inherited link law and $\mathsf H_a'=\mathsf H_a g_i^{-1}$ imply exactly the transformation of $\mathsf A_{aj}$ in \eqref{eq:gaugedata}.

\begin{proposition}[Covariance of the complete coordinate construction]
\label{prop:gauge}
Under \eqref{eq:gaugedata}, the negative ELBO, residual costs, scalar auxiliary factors, information budget, and active parameter $t$ are invariant. Both the unconstrained and constrained Gaussian coordinate solutions transform by pushforward.
\end{proposition}
\begin{proof}
Substitution in \eqref{eq:delta} leaves its residual and trace unchanged. Gaussian state densities acquire the same Jacobian $|\det G|^{-1}$, which cancels in density ratios and KL. Therefore \eqref{eq:freeenergy}, \eqref{eq:betastar}, and \eqref{eq:gammastar} are unchanged. Equations~\eqref{eq:precision} and~\eqref{eq:information} give
\begin{equation}
\mathsf M'=G^{-\top}\mathsf M G^{-1},
\mathsf h'=G^{-\top}\mathsf h,
\label{eq:naturaltransform}
\end{equation}
so $\Sig_*'=G\Sig_*G^\top$ and $\m_*'=G\m_*$. The precision interpolation obeys the same identities. The KL boundary equation is invariant, and its unique active solution is therefore unchanged.
\end{proof}

A change of record coordinates is possible if its measures, noise, and offsets are transformed consistently. With fixed numerical Lebesgue conventions, a density value then acquires the associated Jacobian, and an absolute negative log evidence may change by an additive constant. Proposition~\ref{prop:gauge} states scalar invariance with fixed record coordinates. It also concerns fixed passive re-framing; time-dependent frame changes would require an additional temporal transformation law.

The information projection is not itself a passive gauge transformation. Its reference-relative KL decreases when the constraint is active, and distinct inputs on the same Gaussian exponential ray can reach the same budget boundary. The operation can therefore be many-to-one. Carrying $r$ as geometric data makes the rule covariant; it does not make the active state update information preserving. Holding $\mathsf P=I$ numerically fixed under every coordinate change would instead restrict the symmetry to the appropriate stabilizer.

\subsection{What the recognition family omits}
\label{sec:projectionerrors}

Two restrictions should be distinguished. First, even a full population Gaussian is independent of the auxiliary precision under \eqref{eq:recognition}. In a one-label specialization, let $s_a(x)$ be the record's quadratic residual, and let $\alpha_a=(\nu_a+d_a)/2$. The conditional posterior precision law given $x,o$ has rate $(\nu_a+s_a(x))/2$. At the optimal independent Gamma factor, its rate is the expectation of this conditional rate. Direct subtraction of Gamma log densities gives the additional objective term
\begin{equation}
D_a=\alpha_a\left[\log\bigl(\nu_a+\E_Q s_a\bigr)
-\E_Q\log\bigl(\nu_a+s_a\bigr)\right]\ge0.
\label{eq:gammaerror}
\end{equation}
Indeed, if the conditional rate is $b(x)$ and $\bar b=\E_Q b(x)$, the conditional KL is $\alpha_a\log(\bar b/b(x))-\alpha_a+\alpha_a b(x)/\bar b$; expectation cancels its last two terms. Equation~\eqref{eq:gammaerror} is also the Jensen gap of the logarithm. It vanishes precisely when the residual is constant $Q$-almost surely. A nonzero measurement map and full-support Gaussian give a nonconstant quadratic residual and a strictly positive gap.

Second, restricting $Q$ to a product of site Gaussians discards cross-agent covariance. At fixed source and reliability factors, write the full target precision as $\mathsf M$ with diagonal blocks $\mathsf M_{ii}$. The optimal product Gaussian has the same mean $\mathsf M^{-1}\mathsf h$ and block covariance $\diag(\mathsf M_{ii}^{-1})$. Its minimum excess objective over the unrestricted Gaussian coordinate is
\begin{equation}
\Delta_{\mathrm{site}}=\frac12\log\frac{\prod_i\det\mathsf M_{ii}}{\det\mathsf M}\ge0.
\label{eq:siteerror}
\end{equation}
To derive this expression, substitute the product covariance in $\KL(Q\Vert Q_*)$. The trace becomes $D$, the optimized mean term is zero, and the log-determinant term remains. Nonnegativity follows because the expression is exactly a KL; equality requires the target itself to factor across sites. This comparison concerns the unconstrained Gaussian block at fixed auxiliary factors. It is neither the total posterior error nor the difference between fully reoptimized constrained models. Neither formula authorizes ignoring correlations without an approximation statement.

\section{Gaussian Fisher dynamics, residual corrections, and inference depth}
\label{sec:geometry}

For the full Gaussian family in mean/covariance coordinates, the Fisher metric on tangents $(u,B)$ and $(v,C)$ is
\begin{equation}
g_{(\m,\Sig)}((u,B),(v,C))
=u^\top\Sig^{-1}v+\frac12\Tr(\Sig^{-1}B\Sig^{-1}C).
\label{eq:fisher}
\end{equation}
It follows directly by differentiating the Gaussian log density: the mean score is linear in $x-\m$, the covariance score is a centered quadratic, their cross expectation vanishes by the third centered moment, and the fourth Gaussian moment gives the trace term. This is the Fisher metric in the moment chart, consistent with natural-gradient geometry \citep{Amari1998}. A coordinate gradient in an arbitrary chart is not interchangeable with it.

For a differentiable objective $F$, write $a_F=\partial_{\m}F$ and $G_F=\partial_{\Sig}F$, with the covariance derivative paired by $\Tr(G_F\delta\Sig)$. Metric duality gives
\begin{equation}
\dot\m=-\Sig a_F, \dot\Sig=-2\Sig G_F\Sig,
\label{eq:fisherflow}
\end{equation}
and
\begin{equation}
\frac{dF}{ds}=-a_F^\top\Sig a_F
-2\|\Sig^{1/2}G_F\Sig^{1/2}\|_F^2\le0.
\label{eq:dissipation}
\end{equation}
Here $s$ is inference time. No identification with physical time, posterior sampling, or a stochastic process is made.

For one frozen Gaussian coordinate $F(Q)=\KL(Q\Vert Q_*)+c$, with target precision $\mathsf M$ and information vector $\mathsf h$, direct differentiation gives
\begin{equation}
\dot\m=-\Sig(\mathsf M\m-\mathsf h),
\dot\Sig=\Sig-\Sig\mathsf M\Sig.
\label{eq:frozenflow}
\end{equation}
Introduce $J=\Sig^{-1}$ and $\eta=J\m$. Matrix differentiation then yields the linear equations
\begin{equation}
\dot J=\mathsf M-J, \dot\eta=\mathsf h-\eta.
\label{eq:naturalflow}
\end{equation}
For a duration $s\ge0$, put $\rho=1-e^{-s}$. Their exact solution and mean increment are
\begin{align}
J^+&=(1-\rho)J+\rho\mathsf M,&
\eta^+&=(1-\rho)\eta+\rho\mathsf h,\label{eq:frozenupdate}\\
\m^+-\m&=\rho(J^+)^{-1}(\mathsf h-\mathsf M\m).
\label{eq:residual}
\end{align}
Positive definiteness follows from convexity of the precision cone. The residual is a precision-preconditioned correction toward the coordinate target. It is derived from a frozen objective, not identified with an arbitrary additive message skip. The exact coordinate optimum is approached as $s\to\infty$; finite $s$ is a choice of inference duration. When other factors change during the interval, $\mathsf M$ and $\mathsf h$ change and \eqref{eq:frozenupdate} is not the exact simultaneous solution.

The projection of Theorem~\ref{thm:projection} also has an information-geometric interpretation. Apply \eqref{eq:naturalflow} with target $r$ and initial law $q$. The natural parameters interpolate between $q$ and $r$ with $t=e^{-s}$, precisely giving \eqref{eq:projection}. An active positive budget corresponds to stopping this flow at its unique boundary time. A zero budget is reached only in the infinite-time limit unless the initial law already equals $r$. This calculation does not identify the path with a geodesic of the complete Gaussian Fisher manifold.

Repeated inference reuses the same record laws and recomputes source and reliability factors. Its number of iterations can play the computational role of depth. The negative ELBO does not select a stopping tolerance, maximum compute budget, or synchronous versus asynchronous schedule. Untied depth-dependent record maps generally define different objectives unless a larger model accounts for their variation. The budgeted exact-coordinate theorem also does not follow for an arbitrary truncated flow step that leaves the feasible set.

\section{An independent energy-based construction with a global flow theorem}
\label{sec:globalflow}

A second construction establishes a stronger continuous-time existence statement under different assumptions. It is useful for separating the existence of stable nonlinear agent dynamics from a particular auxiliary-variable representation. Its global-flow conclusion does not automatically apply to the constrained coordinate process above.

Choose a proper product Gaussian anchor $r_G=\mathcal N(a,R_0)$ with block-diagonal $R_0\succ0$. Let $z_{ih}$ range over finite source palettes with positive priors $\pi_{ihj}$. A palette may include a label $0$ with zero interaction. For nonzero sources declare links $L_{ij}^{(h)}$ and receiving metrics $Q_{ihj}\succ0$, and define
\begin{align}
v_{ihj}(x)&=\frac12(x_i-L_{ij}^{(h)}x_j)^\top Q_{ihj}^{-1}(x_i-L_{ij}^{(h)}x_j),\label{eq:edgepotential}\\
U_i(x_i)&=\frac{\chi_i}{4}(x_i^\top\mathsf M_i x_i)^2,\label{eq:quartic}\\
V(x,z)&=\sum_i U_i(x_i)+\sum_{i,h}v_{ih,z_{ih}}(x),
\label{eq:energy}
\end{align}
where $\mathsf M_i\succ0$, $\chi_i\ge0$, and at least one $\chi_i>0$ for the local nonlinear witness. These potentials are specified model assumptions. The normalized joint is
\begin{equation}
p_G(x,z)=Z^{-1}r_G(x)\pi(z)e^{-V(x,z)},
Z=\sum_z\int r_G(x)\pi(z)e^{-V(x,z)}dx.
\label{eq:gibbsjoint}
\end{equation}
Because $V$ is finite and nonnegative, $0<Z\le1$. The graph may have cycles. Absolute polynomial moments are bounded by their Gaussian-anchor moments divided by $Z$. Attach a strictly positive normalized finite-alphabet observation kernel $\ell(o\mid x)$. For the following theorem assume that the resulting Gaussian expectation and its required first and second parameter derivatives admit local integrable domination, giving a $C^2$ objective on $\R^D\times\SPD(D)$. Affine-softmax observation kernels satisfy the needed polynomial-growth conditions; a constant positive kernel gives an explicitly polynomial-moment specialization.

For recognition $Q(x)\prod_{ih}\beta_{ih}(z_{ih})$, let $E_{ihj}=\E_Qv_{ihj}$. Exact minimization of the source rows gives
\begin{equation}
\beta_{ihj}^*=\frac{\pi_{ihj}e^{-E_{ihj}}}{\sum_k\pi_{ihk}e^{-E_{ihk}}}.
\label{eq:gibbssources}
\end{equation}
The reduced negative ELBO is
\begin{equation}
\begin{split}
f(\m,\Sig)={}&\KL(Q\Vert r_G)+\E_Q\left[\sum_iU_i-\log\ell(o\mid x)\right]\\
&-\sum_{ih}\log\sum_j\pi_{ihj}e^{-E_{ihj}}+\log Z.
\end{split}
\label{eq:reduced}
\end{equation}
This follows from the row identity $\sum_j\beta_jE_j+\KL(\beta\Vert\pi)=\KL(\beta\Vert\beta^*)-\log\sum_j\pi_je^{-E_j}$. Every term of \eqref{eq:reduced} except $\log Z$ is nonnegative. In particular, $E_{ihj}\ge0$ implies $\sum_j\pi_{ihj}e^{-E_{ihj}}\le1$.

For $b_i=\m_i^\top\mathsf M_i\m_i+\Tr(\mathsf M_i\Sig_{ii})$, the Gaussian fourth-moment identity gives
\begin{equation}
\E_QU_i=\frac{\chi_i}{4}
\left[b_i^2+2\Tr(\mathsf M_i\Sig_{ii}\mathsf M_i\Sig_{ii})
+4\m_i^\top\mathsf M_i\Sig_{ii}\mathsf M_i\m_i\right].
\label{eq:quarticmoment}
\end{equation}
Its mean derivative is $\chi_i[b_i\mathsf M_i\m_i+2\mathsf M_i\Sig_{ii}\mathsf M_i\m_i]$, containing a cubic term even at fixed covariance. Thus this model also supplies nonlinear local response without inserting a coordinate activation, although here the quartic potential is a declared nonlinear model ingredient.

\begin{theorem}[Global Gaussian Fisher flow]
\label{thm:global}
For fixed model data satisfying the preceding assumptions and any initial $(\m(0),\Sig(0))$ with $\Sig(0)\succ0$, the Fisher flow \eqref{eq:fisherflow} for $f$ has a unique solution for every $s\ge0$. It stays in a compact subset of $\R^D\times\SPD(D)$, decreases $f$, and has only stationary accumulation points.
\end{theorem}
\begin{proof}
The nonnegative terms in \eqref{eq:reduced} imply
\begin{equation}
f(\m,\Sig)\ge\KL(Q\Vert r_G)+\log Z.
\label{eq:coercivity}
\end{equation}
The dissipation identity \eqref{eq:dissipation} therefore confines the trajectory to a Gaussian-KL sublevel. The spectral argument of \eqref{eq:spectralbound} bounds the mean and covariance above and keeps the covariance away from zero. The closed objective sublevel is a compact subset of the open parameter domain.

The $C^2$ assumption makes the vector field locally Lipschitz. Local existence and uniqueness follow by contraction of its integral equation on a sufficiently small parameter ball. On the compact sublevel, the velocity is bounded. If a maximal solution ended at finite time, it would have an interior limit and the same local construction would extend it, a contradiction. This proves global existence and positive-definite covariance for all finite times.

Let $\mathcal D(s)=-df/ds\ge0$ be the right-hand dissipation in \eqref{eq:dissipation}. The lower bound on $f$ implies $\int_0^\infty\mathcal D(s)ds<\infty$. The dissipation and vector field are continuously differentiable with bounded derivatives on the compact sublevel, so $\mathcal D(s)$ is uniformly continuous. If it did not tend to zero, uniform continuity would give disjoint intervals carrying a fixed positive amount of dissipation, contradicting integrability. Hence $\mathcal D(s)\to0$. At an accumulation point, positive-definite covariance makes both objective derivatives zero.
\end{proof}

The theorem asserts a unique trajectory for each initial condition, not a unique equilibrium or a global minimizer. It does not establish convergence to one point. Under fixed passive changes, transform the anchor by pushforward, $\mathsf M_i'=g_i^{-\top}\mathsf M_i g_i^{-1}$, $Q_{ihj}'=g_iQ_{ihj}g_i^\top$, and links by their endpoint law. Then $V'(Gx,z)=V(x,z)$ and $Z'=Z$. Transforming the observation kernel consistently gives an invariant $f$. Its mean and covariance derivatives transform by the corresponding inverse actions, and substituting in \eqref{eq:fisherflow} proves covariance of the deterministic trajectories.

\subsection{Finite steps and their computational assumptions}

At a nonstationary point, with $a_f=\partial_{\m}f$ and $G_f=\partial_{\Sig}f$, consider
\begin{equation}
\m_\eta=\m-\eta\Sig a_f,
\Sig_\eta=\Sig^{1/2}\exp(-2\eta\Sig^{1/2}G_f\Sig^{1/2})\Sig^{1/2}.
\label{eq:retraction}
\end{equation}
The covariance is positive definite for every finite $\eta\ge0$, and the derivative at zero is the Fisher vector field. The identity $\Sig_\eta=\Sig\exp(-2\eta G_f\Sig)$ and matrix similarity prove covariance under passive frame changes. This is a covariance retraction combined with a mean increment, not the exact exponential map of the entire Gaussian Fisher manifold.

For fixed $0<\theta<1$, accept a step only if
\begin{equation}
f(\m_\eta,\Sig_\eta)\le f(\m,\Sig)-\theta\eta\mathcal D.
\label{eq:armijo}
\end{equation}
Its derivative at zero is $-\mathcal D<0$, so geometric backtracking terminates under exact evaluation. A certified numerical version requires error bounds sufficient to validate the acceptance inequality. General Gaussian expectations of log-softmax likelihoods have no finite exact evaluation rule supplied here. The constant-observation specialization avoids that particular integral because the remaining expectations are polynomial moments and finite source sums; the fixed additive $\log Z$ cancels in objective differences. Matrix-function and floating-point errors still require numerical control.

Positive-definite covariance alone does not ensure descent. In one variance coordinate, take $f(c)=(c-1-\log c)/2$, whose Fisher flow is $\dot c=c-c^2$. From $c=2$, an Euler step of size $2$ gives $c^+=-2$. The exponential covariance step instead gives $c^+=2e^{-2}>0$, but its objective increment is exactly $e^{-2}>0$. Both witnesses embed into a larger population by leaving other coordinates unchanged.

A constrained proximal prescription provides another analytical descent statement. On the compact set $\mathcal G_B=\{Q\text{ Gaussian}:\KL(Q\Vert r_G)\le B\}$, continuity gives a minimizer of $f(Q)+\eta^{-1}\KL(Q\Vert Q_n)$. Comparing it with the feasible competitor $Q_n$ proves
\begin{equation}
f(Q_{n+1})+\eta^{-1}\KL(Q_{n+1}\Vert Q_n)\le f(Q_n).
\label{eq:proximal}
\end{equation}
This is an existence result; a global nonlinear argmin is not a supplied efficient solver. If minimizers are nonunique, the solution set transforms covariantly, while a single-valued selection requires an additional rule. This distinction prevents an optimization oracle from being mistaken for a derived computational architecture.

\section{A categorical observation model without an MLP decoder}
\label{sec:readout}

A language model needs a probability law on discrete observations. For a finite vocabulary of size $V\ge2$, attach utilities
\begin{equation}
u\mid x\sim\mathcal N(Hx+c,S), S\in\SPD(V),
 y=\operatorname*{argmax}_{1\le b\le V}u_b.
\label{eq:utility}
\end{equation}
The matrix $H:\R^D\to\R^V$ and vector $c$ are declared observation parameters. Let $\mathcal K_y=\{u:u_y>u_b\text{ for all }b\ne y\}$. Every difference $u_y-u_b$ has positive variance, so ties have probability zero. The cones are disjoint and cover the utility space except for null boundaries. Each contains an open ball of positive Gaussian probability. Thus their probabilities define a strictly positive normalized categorical law. The latent-Gaussian augmentation follows a classical statistical route \citep{AlbertChib1993}; the utility coordinates are observation variables, not a prescribed hidden neuron layer.

Given an observed category, extend recognition by an independent density $\xi(u)$ supported on $\mathcal K_y$, with finite entropy and second moments. At fixed $Q$, the expected quadratic satisfies
\begin{equation}
\E_Q[(u-Hx-c)^\top S^{-1}(u-Hx-c)]
=(u-H\m-c)^\top S^{-1}(u-H\m-c)+\Tr(S^{-1}H\Sig H^\top).
\label{eq:utilityquadratic}
\end{equation}
The unique utility coordinate is therefore
\begin{equation}
\xi_*(u)=\frac{1_{\mathcal K_y}(u)\mathcal N(u;H\m+c,S)}
 {\Pr_{\mathcal N(H\m+c,S)}(\mathcal K_y)}.
\label{eq:truncated}
\end{equation}
The normalizer is positive; Gaussian tail bounds give finite moments, and its quadratic log density on the cone gives finite entropy.

To make the ELBO accounting explicit, write $\bar u=\E_\xi u$ and $W=\Cov_\xi(u)$. The additional contribution is
\begin{equation}
\begin{split}
\mathcal F_u={}&\E_\xi\log\xi+\frac12\bigl[V\log(2\pi)+\log|S|+\Tr(S^{-1}W)\\
&+(\bar u-H\m-c)^\top S^{-1}(\bar u-H\m-c)
 +\Tr(S^{-1}H\Sig H^\top)\bigr].
\end{split}
\label{eq:utilityelbo}
\end{equation}
It is infinite when $\xi$ gives positive mass outside the observed cone. Completing the density ratio rewrites this as $\KL(\xi\Vert\xi_*)-\log\Pr_{\mathcal N(H\m+c,S)}(\mathcal K_y)+\Tr(S^{-1}H\Sig H^\top)/2$. The truncation normalizer is already accounted for through the optimized factor's entropy. Adding it again to the uncollapsed joint would double-count it.

At fixed $\xi$, the state-coordinate precision and information vector receive the additions
\begin{equation}
\mathsf M\longmapsto\mathsf M+H^\top S^{-1}H,
\mathsf h\longmapsto\mathsf h+H^\top S^{-1}(\bar u-c).
\label{eq:utilitystate}
\end{equation}
The precision stays positive definite, the coordinate remains Gaussian, and Theorem~\ref{thm:projection} applies without changing its proof. Consequently the observation extension retains exact sequential common-objective descent, assuming exact utility moments. Passive covariance uses $H'=HG^{-1}$, with fixed utility coordinates and $S,c$.

For prediction of a fresh unobserved utility, integrate it under $Q$:
\begin{equation}
u\sim\mathcal N(H\m+c,S+H\Sig H^\top),
\widehat p(y\mid Q)=\Pr_{\mathcal N(H\m+c,S+H\Sig H^\top)}(\mathcal K_y).
\label{eq:utilitypredictive}
\end{equation}
This is a normalized posterior-predictive approximation. Its covariance differs from the covariance $S$ in the auxiliary coordinate \eqref{eq:truncated}: prediction averages a density, while variational coordinate optimization exponentiates an expected log density. For two categories the probability is an ordinary univariate Gaussian CDF,
\begin{equation}
\widehat p(1\mid Q)=\Phi\left(
\frac{(h_1-h_2)\m+c_1-c_2}
{\sqrt{(e_1-e_2)^\top(S+H\Sig H^\top)(e_1-e_2)}}\right),
\label{eq:binaryprobit}
\end{equation}
where $h_b$ is row $b$ of $H$ and $e_b$ is a coordinate vector. No softmax equality is asserted. For a large vocabulary, Gaussian cone probabilities and moments are a substantial computational obligation. Vocabulary semantics, parameter identifiability, and the choice of observation law remain unresolved by gauge symmetry.
\section{Causal prediction and the sequence-model obligation}
\label{sec:causality}

A triangular source graph restricts which links enter receiver costs. It does not make a complete population inference gradient causal. For the allowed edge $2\leftarrow1$, the quadratic cost
\begin{equation}
E_{21}=\frac12(\m_2-L_{21}\m_1)^\top R^{-1}(\m_2-L_{21}\m_1)
\label{eq:causalcost}
\end{equation}
has source-side derivative
\begin{equation}
\partial_{\m_1}E_{21}=-L_{21}^\top R^{-1}(\m_2-L_{21}\m_1).
\label{eq:sourcegradient}
\end{equation}
If agent $2$ has assimilated a future observation, the complete gradient changes the earlier agent using that information. Dropping source derivatives may define a different directed algorithm, but it would no longer be the full gradient of the same objective. Correlated population covariance can introduce additional coupling in the Fisher update.

A sufficient probabilistic construction is to compute each predictive state $Q_t$ using only the available prefix $y_{<t}$ and then evaluate the same state for every candidate next category. A declared deterministic inference algorithm gives a normalized kernel $k_t(y_t\mid y_{<t})$ through \eqref{eq:utilitypredictive}. For fixed finite sequence length $T$,
\begin{equation}
\widehat p(y_{1:T})=\prod_{t=1}^{T}k_t(y_t\mid y_{<t})
\label{eq:chain}
\end{equation}
is normalized by successively summing $y_T,y_{T-1},\ldots,y_1$. A randomized inference algorithm requires a specified seed convention or an explicit marginal over its randomness. A law on finite strings of variable length additionally requires an almost-sure stopping mechanism.

The prefix requirement is substantive. For a one-token binary problem, let inference inspect the candidate target and assign it probability $3/4$. Each separately computed categorical row is normalized, but the resulting scores for the two possible tokens sum to $3/2$. A target-dependent state is not a single predictive kernel conditioned on the empty prefix.

Equation~\eqref{eq:chain} establishes a sequence distribution. It does not establish that this approximate predictive chain equals the exact sequence marginal of \eqref{eq:jointlaw}, nor that its negative log likelihood equals a sum of the training ELBOs above. A language-model construction must specify which records encode the prefix, how new states and links are introduced, and which conditional joint is bounded during learning. Efficient online reuse of one inference state across prefixes also needs its own covariance and objective argument. Re-solving a full population problem at each prefix proves an existence route for predictions, not practical transformer-scale computation.

\section{Why standard blocks are not universal identities}
\label{sec:obstructions}

\subsection{LayerNorm, RMSNorm, and full Gaussian laws}

LayerNorm \citep{Ba2016} centers a vector across its components before rescaling. With zero gain/bias modifications suppressed, define $P_0=I-\mathbf1\mathbf1^\top/K$ and
\begin{equation}
\operatorname{LN}_{\epsilon}(x)=\frac{P_0x}
 {\sqrt{\|P_0x\|^2/K+\epsilon}}.
\label{eq:layernorm}
\end{equation}
For $K=2$ and $\epsilon=0$, a nondegenerate Gaussian has $x_1\ne x_2$ almost surely and
\begin{equation}
\operatorname{LN}_{0}(x)=\operatorname{sign}(x_1-x_2)(1,-1).
\label{eq:lntwopoint}
\end{equation}
Its output has two atoms, not a full Gaussian law. With $\epsilon>0$, the output still lies in the proper zero-sum hyperplane. Fixed affine gains and biases map that support into an affine subspace of dimension at most $K-1$, so they cannot restore nondegenerate full Gaussian support.

RMSNorm omits centering \citep{ZhangSennrich2019}. At $\epsilon=0$, its normalized output lies on a sphere. At $\epsilon>0$, its norm is strictly below $\sqrt K$. The map can then be injective: bounded support alone is not proof of information loss. In either case, however, its output is not a full nondegenerate Gaussian law. These are obstructions to exact pushforward closure, not to defining a Gaussian approximation after normalization or using a coordinate operation on the recognition mean by convention.

The reference-budget rule supplies a different statistical operation with a proved scale-control property. Its invariance follows from transforming the reference and the entire belief law. It does not reproduce componentwise centering, affine gains, or the optimization behavior of LayerNorm, and no equality of training benefit has been shown.

\subsection{Arbitrary MLP fields and scalar-potential integrability}

Suppose a differentiable mean field $f(\mu)$ is asserted to equal $-\Sigma\nabla\Phi(\mu)$ for a scalar potential under a fixed SPD Gaussian mean metric. Then $\Sigma^{-1}Df$ must be symmetric. Consider the realizable one-direction field
\begin{equation}
f(\mu_1,\mu_2)=(0,\sigma(\mu_1)),
\Sigma^{-1}=\begin{pmatrix}a&b\\b&d\end{pmatrix}\succ0.
\label{eq:mlpfield}
\end{equation}
The off-diagonal entries of $\Sigma^{-1}Df$ are $0$ and $d\sigma'(\mu_1)$. Since $d>0$, they disagree wherever $\sigma'\ne0$. Such an MLP field cannot be the claimed fixed-metric scalar-potential gradient. This excludes a universal identification of arbitrary FFNs with one such inference step; it does not exclude augmented latent models, external readouts, or other metric-dependent constructions.

There is a further full-state obstruction to a nonzero covariance-independent mean field with universal covariance passthrough. Under \eqref{eq:fisherflow}, $\dot\Sigma=0$ for every SPD $\Sigma$ implies $\partial_\Sigma F=0$. Connectedness of the SPD cone then makes $F$ depend only on the mean. Requiring $f(\mu)=-\Sigma\nabla_\mu F$ to be the same at both $\Sigma$ and $2\Sigma$ forces $f=0$. This statement concerns an instantaneous vector field on the full Gaussian manifold. It is not a theorem about every composite finite endpoint map.

\subsection{Scope of the earlier gated-residual analogy}

The prior gauge-VFE manuscript \path{PIFB2.tex}, equation labeled \path{eq:vfe_glu}, already relates Boltzmann-gated prediction errors to a GLU-like computational pattern. That relationship is a useful precursor to the present investigation. A generic assertion that each binary-source gated residual is odd and redescending in both tails needs correction. For scalar sources at $0$ and $1$ with equal variances and temperature one, the source-$0$ contribution is
\begin{equation}
f(x)=x\frac{e^{-x^2/2}}{e^{-x^2/2}+e^{-(x-1)^2/2}}
=\frac{x}{1+e^{x-1/2}}.
\label{eq:gatecounter}
\end{equation}
It is not odd, and $f(x)/x\to1$ in the negative tail. The positive tail vanishes. This counterexample limits a property of the per-edge expression; it does not remove the valid gating mechanism or its explicitly qualified GLU-family analogy. The symmetric Student-t response in \eqref{eq:scalarresponse} is a different mechanism.

A derivative distinction also matters. For fixed source prior and $\beta_j=\pi_je^{-E_j/\tau}/\sum_k\pi_ke^{-E_k/\tau}$, put $\bar E=\sum_j\beta_jE_j$. Differentiating the entropy-free surrogate gives
\begin{equation}
\nabla\sum_j\beta_jE_j
=\sum_j\beta_j\left[1-\frac{E_j-\bar E}{\tau}\right]\nabla E_j.
\label{eq:surrogate}
\end{equation}
To obtain this identity, substitute $\nabla\beta_j=-(\beta_j/\tau)(\nabla E_j-\sum_k\beta_k\nabla E_k)$. The simpler $\sum_j\beta_j\nabla E_j$ is the derivative with fixed responsibilities or the envelope derivative of the optimized entropy-augmented objective. The companion work already distinguishes these routes; they must remain distinct when interpreting a nonlinear update.

\section{Architectural consequences and remaining choices}
\label{sec:choices}

The central result supplies one normalized law whose exact sequential inference performs source selection, nonlinear reliability processing, covariance revision, and reference-relative scale control. The specified generative state-to-record maps are affine; the nonlinearity is derived from the declared latent-variable inference. Thus prescribed neuron layers are unnecessary for existence of these roles. This conclusion is narrower than deriving an entire language-model architecture from gauge invariance or showing that the current coordinate MLP is dispensable in its empirical configuration.

\begin{table}[tbp]
\centering
\small
\renewcommand{\arraystretch}{1.15}
\begin{tabularx}{\textwidth}{@{}>{\raggedright\arraybackslash}p{0.15\textwidth}>{\raggedright\arraybackslash}X>{\raggedright\arraybackslash}X@{}}
\toprule
Component & Conditional agent-theory mechanism & Remaining choice or limitation \\
\midrule
Attention & Source-label coordinate probabilities. & Palette, prior, and record cost; generally different from marginal Gaussian KL. \\
MLP role & Nonlinear precision and source inference with covariance revision. & Expressive sufficiency and statistical selection of the latent law. \\
Normalization & Exact reference-KL constrained Gaussian coordinate. & Reference, budget, and divergence orientation; no LayerNorm identity. \\
Residual & Preconditioned correction along a specified inference dynamics. & Duration and accepted discretization; no arbitrary skip equivalence. \\
Depth & Repetition under a fixed objective. & Stopping criterion, scheduling, and computational budget. \\
Heads & Multiple typed record and source channels. & Number, rank, overlap, and parameter sharing. \\
Positions & Declared order-dependent data, links, or source priors. & No selection theorem for RoPE or a particular positional law. \\
Causality & Prefix-only inference with one state shared across candidate outputs. & Efficient online reuse; source masks alone are insufficient. \\
Output & Gaussian utility regions with normalized probabilities. & Vocabulary semantics, identifiability, and cone integration cost. \\
Learning & Optimization of a declared normalized model and recognition family. & Sequence likelihood relation, scalable solver, and generalization. \\
\bottomrule
\end{tabularx}
\caption{Functional roles supplied by the conditional constructions and the choices that remain. Rows are not claims of exact equality with the named transformer blocks. The global-flow theorem belongs to its separate energy-based model.}
\label{tab:roles}
\end{table}

Table~\ref{tab:roles} separates the supplied mechanisms from the remaining choices. The starting principles underdetermine even the reliability response. Keep the same reference, maps, receiving scale $R$, and passive transformation rules, but choose $\nu=4$ or $\nu=8$. Both models are normalized and admit the same full Gaussian recognition domain. For Student noise $\varepsilon=R^{1/2}Z/\sqrt\lambda$, with $Z\sim\mathcal N(0,I)$ independent of $\lambda$, Gamma integration gives
\begin{equation}
\Cov(\varepsilon)=\E(\lambda^{-1})R=\frac{\nu}{\nu-2}R, \nu>2.
\label{eq:underdetermined}
\end{equation}
The two conditional observation laws have covariances $2R$ and $4R/3$, hence differ. Their scalar one-sweep means at $y=1$ are $5/11$ and $9/19$. The interacting contrast embeds both responses into the allowed population dimensions. Therefore gauge covariance, normalization, and the Gaussian recognition domain do not uniquely select this local inference rule. Additional observation semantics or selection principles could distinguish them; this witness does not rule out such further theory.

Several computational costs remain visible. Full population covariance costs more than storing independent site covariances; a sparse approximation needs an error statement. Exact auxiliary coordinates use current factors; parallel stale updates require a separate descent analysis. The information constraint introduces a scalar root problem, the categorical extension introduces Gaussian cone moments, and a generic nonlinear proximal prescription introduces a global optimization problem. An existence theorem does not remove these costs by writing them as named operators.

Parameter learning must use the stated normalized law, including every parameter-dependent normalization and entropy term. In the energy model, changing the potential changes $Z$; its derivative cannot be omitted simply because $Z$ is constant during fixed-model inference. In the utility model, one must differentiate either the full joint objective with appropriate auxiliary factors or the correctly profiled objective, retaining the probability and trace terms from \eqref{eq:utilityelbo}. No optimizer or global convergence theorem for these learning problems is supplied here.

The next mathematical task is a fully specified prefix model with a declared observation ontology, stored covariance structure, finite evaluation rule, and learning objective. Its sequence likelihood must be related explicitly to its training ELBO and predictive approximation. Task-level usefulness then requires controlled empirical comparisons. The existing transformer development results belong to a different implementation containing coordinate MLPs and covariance passthrough; they are not evidence for the performance of the model derived in this extension.

\section{Summary of the conditional result}

For every finite $N\ge2$ and $K\ge2$, the normalized construction admits a nonlinear interacting full-Gaussian recognition response with a finite reference-KL budget. Source, reliability, and Gaussian factors have exact compatible coordinate updates, and the information constraint is solved within that same objective. Passive local coordinate changes preserve the scalar objective and transform the full belief law covariantly. A separate normalized energy model gives global Gaussian Fisher trajectories and stationary accumulation points under its own regularity and coercivity assumptions. A latent-Gaussian utility extension supplies normalized categorical observations while leaving their efficient evaluation open.

These results extend the gauge-VFE account beyond a geometric attention operator to conditional population inference mechanisms. They do not supply an exact derivation of the companion implementation's LayerNorm or arbitrary MLP, a unique architectural selection principle, or an efficient competitive language model. The distinction between mathematical compatibility, exact coordinate inference, numerical realization, and empirical adequacy is part of the construction's scope.

\appendix
\section{Finite checks and reproducibility}
\label{app:checks}

The derivations above were accompanied by small CPU calculations. These calculations check formulas and representative implementations; they do not replace the quantified proofs or constitute language-model experiments. A three-agent example with two state coordinates per agent and two source labels per record executed $90$ sequential coordinate replacements with information budget $B=0.24$. The negative ELBO decreased from $10.7434417243$ to $9.01470261093$. Its largest measured increase was $1.78\times10^{-15}$, consistent with floating-point roundoff. The largest observed budget value was $0.24000000000000066$, so this numerical root solve was only accurate to floating-point tolerance, rather than an exact feasibility certificate.

The final covariance had minimum eigenvalue $0.3467039505$ and a nonzero cross-agent block of norm $0.1029894600$. Passive-frame discrepancies in means, covariance, routing, reliability, and objective were at most $3.56\times10^{-15}$. Low-dimensional quadrature checked the Gaussian--Gamma Student mixture, a truncated utility mean, and the binary predictive probability with errors at most $1.53\times10^{-16}$ for those particular inputs. A separate exact rational-matrix calculation checked ten normalization, covariance, and Fisher-coordinate identities. The local reliability calculation also checked the fixed-weight product-family gap in \eqref{eq:siteerror}.

The V3\_Transformer repository retains scripts, machine-readable outputs, source snapshots, and independent derivations in its dated research package, \path{docs/derivations/2026-09-20-neuron-free-agent-theory/}. Within that directory, the joint-model record is \path{evidence/hybrid-check-results.json}. The separate exact normalization and local reliability records are \path{agents/normalization-check-results.json} and \path{agents/nonlinear-evidence.json}, respectively. The research report records the full architectural target as unresolved while identifying the conditional theorems proved here. No training, GPU experiment, or Lean formalization was performed for this extension. The present manuscript does not alter the existing implementation or its configuration.

\section{Why finite auxiliary initialization is necessary}

The integrability condition preceding \eqref{eq:delta} cannot be replaced by Gaussian budget feasibility. For example, take a feasible $Q=r$, nonzero record maps, and an auxiliary half-Cauchy precision density $\gamma(\lambda)=2/[\pi(1+\lambda^2)]$ on $\lambda>0$. It has infinite $\E\lambda$ while the expected quadratic residual is positive. Every exponent in the naive source-update numerator is then negative infinity, producing an undefined zero-over-zero normalization. Its KL to the Gamma prior is also infinite. Requiring finite initial free energy and the stated moments excludes this case. Positive finite-shape/rate Gamma initialization gives a concrete admissible choice, and \eqref{eq:gammastar} propagates those conditions through every finite exact sweep.

\begin{thebibliography}{99}
\raggedright
\bibitem[Albert and Chib(1993)]{AlbertChib1993}
James H. Albert and Siddhartha Chib.
Bayesian analysis of binary and polychotomous response data.
\emph{Journal of the American Statistical Association}, 88(422):669--679, 1993.
\href{https://stat.cmu.edu/~brian/905-2009/all-papers/albert-chib-1993.pdf}{Primary article}.

\bibitem[Amari(1998)]{Amari1998}
Shun-ichi Amari.
Natural gradient works efficiently in learning.
\emph{Neural Computation}, 10(2):251--276, 1998.
\href{https://doi.org/10.1162/089976698300017746}{doi:10.1162/089976698300017746}.

\bibitem[Andrews and Mallows(1974)]{AndrewsMallows1974}
D. F. Andrews and C. L. Mallows.
Scale mixtures of normal distributions.
\emph{Journal of the Royal Statistical Society, Series B}, 36(1):99--102, 1974.
\href{https://doi.org/10.1111/j.2517-6161.1974.tb00989.x}{doi:10.1111/j.2517-6161.1974.tb00989.x}.

\bibitem[Ba et~al.(2016)Ba, Kiros, and Hinton]{Ba2016}
Jimmy Lei Ba, Jamie Ryan Kiros, and Geoffrey E. Hinton.
Layer normalization.
arXiv:1607.06450, 2016.
\url{https://arxiv.org/abs/1607.06450}.

\bibitem[Blei et~al.(2017)Blei, Kucukelbir, and McAuliffe]{Blei2017}
David M. Blei, Alp Kucukelbir, and Jon D. McAuliffe.
Variational inference: A review for statisticians.
\emph{Journal of the American Statistical Association}, 112(518):859--877, 2017.
\href{https://doi.org/10.1080/01621459.2017.1285773}{doi:10.1080/01621459.2017.1285773}.

\bibitem[Wetterich(2014)]{Wetterich2014}
Christof Wetterich.
Linear lattice gauge theory.
\emph{Nuclear Physics B}, 884:44--65, 2014.
\href{https://doi.org/10.1016/j.nuclphysb.2014.04.002}{doi:10.1016/j.nuclphysb.2014.04.002}.

\bibitem[Yu et~al.(2024)Yu, Buchanan, Pai, Chu, Wu, Tong, Bai, Zhai, Haeffele, and Ma]{Yu2024CRATE}
Yaodong Yu, Sam Buchanan, Druv Pai, Tianzhe Chu, Ziyang Wu, Shengbang Tong,
Hao Bai, Yuexiang Zhai, Benjamin D. Haeffele, and Yi Ma.
White-box transformers via sparse rate reduction: Compression is all there is?
\emph{Journal of Machine Learning Research}, 25(300):1--128, 2024.
\url{https://jmlr.org/papers/v25/23-1547.html}.

\bibitem[Zhang and Sennrich(2019)]{ZhangSennrich2019}
Biao Zhang and Rico Sennrich.
Root mean square layer normalization.
\emph{Advances in Neural Information Processing Systems}, 32, 2019.
\url{https://arxiv.org/abs/1910.07467}.
\end{thebibliography}
\end{document}